\chapter*{Project Repositories and Structure}

The source artifacts for this project are split across two GitHub repositories.
The main repository contains the complete benchmark and dissertation workspace,
while the standalone repository packages the final authenticated AES target as a
smaller release artifact.

\begin{tabularx}{\textwidth}{@{}P{1.45in}L@{}}
\textbf{Repository} & \textbf{GitHub link and purpose} \\
\hline
\texttt{fyp} & \url{https://github.com/kush-10/fyp} \\
& Main project repository containing the RISC Zero benchmark workspaces, Python
benchmark harness, generated report artifacts, and dissertation source. \\
\texttt{aes-ctr-hmac-r0} & \url{https://github.com/kush-10/aes-ctr-hmac-r0} \\
& Standalone release repository for the final AES-192-CTR plus
HMAC-SHA256-192 benchmark target. \\
\end{tabularx}

At the top level, \texttt{fyp} is organised around benchmark workspaces such as
\texttt{aes-r0}, \texttt{aes-r0-optimised}, \texttt{aes-ctr},
\texttt{aes-ctr-hmac}, \texttt{lowmc-r0}, \texttt{lowmc-r0-optimised},
\texttt{salsa-r0}, and \texttt{operation-bnchmrk-r0}. Each workspace follows the
same host/guest pattern: the host prepares inputs, runs the RISC Zero prover,
verifies the receipt, and emits timing or JSON benchmark output; the guest
contains the computation proven inside the zkVM.

Shared experiment infrastructure lives in \texttt{bench-harness}, which provides
the config-driven runner, aggregation scripts, plotting scripts, and benchmark
schema notes. Dissertation sources live in \texttt{docs}, and selected benchmark
outputs live below \texttt{artifacts}. The standalone \texttt{aes-ctr-hmac-r0}
repository keeps the same core layout for the final target: \texttt{aesencryption}
for AES/HMAC code, \texttt{host} for proving and verification, \texttt{methods}
for the RISC Zero method wrapper and guest, and \texttt{scripts} for the
standalone report benchmark workflow.
